\chapter{Results and Discussion}

\section{Stage 1 Regression Results}

The Stage 1 model was evaluated as a continuous Fire Weather Index (FWI)
regression model. Its task is to estimate daily FWI from the locked set of 34
engineered input features described in Chapter 3. The deployed model is a
\texttt{HistGradientBoostingRegressor}; its output,
\texttt{predicted\_fwi}, is the numerical risk estimate shown in the dashboard
and passed to the high-risk decision layer.

Table~\ref{tab:stage1-regression-metrics} reports the verified Stage 1
metrics from the repository metadata. The out-of-fold validation result was
RMSE \(=4.111\), MAE \(=2.918\), and \(R^2=0.835\). The 2025 holdout result was
RMSE \(=7.259\), MAE \(=5.098\), and \(R^2=0.819\). The increase in error on
the later holdout period is expected because it tests the model on a
chronologically later season, but the holdout \(R^2\) still indicates that the
model captures much of the variation in the FWI target.

\begin{table}[htbp]
\centering
\caption{Stage 1 regression performance for continuous FWI prediction.}
\label{tab:stage1-regression-metrics}
\begin{tabularx}{\textwidth}{l c c c X}
\toprule
Evaluation setting & RMSE & MAE & \(R^2\) & Interpretation \\
\midrule
Out-of-fold validation & 4.111 & 2.918 & 0.835 &
Temporal validation result across training folds; shows strong fit while
respecting chronological evaluation. \\
2025 holdout test & 7.259 & 5.098 & 0.819 &
Later-season holdout result; error increases, but the model still captures
most variation in the FWI target. \\
\bottomrule
\end{tabularx}
\end{table}


\begin{figure}[htbp]
  \centering
  \includegraphics[width=0.86\textwidth]{stage1_actual_vs_predicted.png}
  \caption{Stage 1 actual versus predicted FWI comparison for candidate regression models.}
  \label{fig:stage1-actual-vs-predicted}
\end{figure}

Figure~\ref{fig:stage1-actual-vs-predicted} shows that the selected regression
approach follows the general trend of observed FWI values. The result should
not be interpreted as exact fire-danger reproduction; instead, it is a
decision-support estimate that preserves useful information about how close a
day is to the high-risk boundary.

\begin{figure}[htbp]
  \centering
  \includegraphics[width=0.86\textwidth]{stage1_residuals.png}
  \caption{Stage 1 residual analysis showing prediction error behaviour across FWI values.}
  \label{fig:stage1-residuals}
\end{figure}

The residual plot in Figure~\ref{fig:stage1-residuals} is important because
wildfire risk decisions are sensitive to larger errors near high FWI values.
The model is useful for daily risk estimation, but the residual pattern also
supports the need for a second safety-oriented classification layer.

\begin{figure}[htbp]
  \centering
  \includegraphics[width=0.82\textwidth]{stage1_regression_metric_comparison.png}
  \caption{Regression benchmark comparison using \(R^2\), RMSE, and MAE.}
  \label{fig:stage1-metrics}
\end{figure}

Figure~\ref{fig:stage1-metrics} summarizes the Stage 1 benchmark evidence. The
combined use of RMSE, MAE, and \(R^2\) is necessary because a single metric
cannot describe both average error and explained variation.

\FloatBarrier

\section{Stage 2 High-Risk Classification Results}

Stage 2 was evaluated as a high-risk classification layer. It uses a
\texttt{RandomForestClassifier} with four inputs:
\texttt{predicted\_fwi}, \texttt{rh}, \texttt{ws}, and
\texttt{fuel\_drying\_rate}. The class target is \texttt{target\_ge\_35},
where a day is high risk when \( \mathrm{FWI} \geq 35 \). The final stored
probability threshold is 0.10.

Table~\ref{tab:stage2-classification-metrics} gives the final classifier
values verified from the repository metadata. At threshold 0.10, the Stage 2
classifier achieved precision \(=0.849\), recall \(=0.900\), F1-score
\(=0.874\), and accuracy \(=0.929\), with 45 true positives, 8 false
positives, 5 false negatives, and 126 true negatives.

\begin{table}[htbp]
\centering
\caption{Final Stage 2 high-risk classification result at probability threshold 0.10.}
\label{tab:stage2-classification-metrics}
\begin{tabularx}{\textwidth}{c c c c c c c c c}
\toprule
Threshold & Precision & Recall & F1-score & Accuracy & TP & FP & FN & TN \\
\midrule
0.10 & 0.849 & 0.900 & 0.874 & 0.929 & 45 & 8 & 5 & 126 \\
\bottomrule
\end{tabularx}
\end{table}


\begin{figure}[htbp]
  \centering
  \includegraphics[width=0.86\textwidth]{stage2_precision_recall_comparison.png}
  \caption{Stage 2 precision--recall comparison for imbalanced high-risk classification.}
  \label{fig:stage2-pr-curve}
\end{figure}

Figure~\ref{fig:stage2-pr-curve} is more informative than accuracy alone
because high-risk days are less frequent than normal days. In this application,
recall is especially important: a false negative may reduce attention on a
genuinely dangerous day. Precision still matters because too many false alarms
can reduce operator trust and increase monitoring load.

\begin{figure}[htbp]
  \centering
  \includegraphics[width=0.68\textwidth]{stage2_final_confusion_matrix_threshold_010.png}
  \caption{Final Stage 2 high-risk classifier confusion matrix using the report threshold of 0.10.}
  \label{fig:stage2-confusion-matrix}
\end{figure}

The confusion matrix in Figure~\ref{fig:stage2-confusion-matrix} uses the
final report values, not the notebook ablation figure at threshold 0.25. This
distinction is important because mixing ablation results with the final
deployed threshold would misrepresent the classifier performance.

\begin{figure}[htbp]
  \centering
  \includegraphics[width=0.72\textwidth]{stage2_final_classifier_metrics.png}
  \caption{Final Stage 2 classifier metrics at threshold 0.10.}
  \label{fig:stage2-classifier-metrics}
\end{figure}

\begin{figure}[htbp]
  \centering
  \includegraphics[width=0.86\textwidth]{feature_importance_ablation.png}
  \caption{Feature importance evidence for the enriched Random Forest classifier, highlighting the contribution of the support features.}
  \label{fig:feature-importance}
\end{figure}

Figures~\ref{fig:stage2-classifier-metrics} and
\ref{fig:feature-importance} show that the final decision layer is not a
black-box replacement for the regression model. It is a compact safety layer
that uses the Stage 1 output plus humidity, wind, and fuel-drying information
to reduce missed high-risk cases near the decision boundary.

\FloatBarrier

\section{YOLOv8 Fire and Smoke Detection Results}

The visual monitoring stage uses the active YOLOv8 \texttt{best.pt} detector
for two object classes: fire and smoke. The detector provides visual evidence
through bounding boxes, class labels, and confidence values. This is different
from the prediction model: prediction estimates fire-weather risk, while YOLOv8
detection identifies visible evidence in camera or drone-ready frames.

Object detection should be evaluated with object-detection metrics rather than
generic image-classification accuracy. Precision measures how many detections
are correct, recall measures how many relevant objects are found, and
mAP@0.5 / mAP@0.5:0.95 summarize localisation and classification quality
across confidence and overlap thresholds. Table~\ref{tab:yolo-detection-metrics}
therefore reports precision, recall, and mAP values from the uploaded YOLOv8
training/evaluation output.

\begin{table}[htbp]
\centering
\caption{YOLOv8 fire/smoke detection metrics from the uploaded training and evaluation output.}
\label{tab:yolo-detection-metrics}
\begin{tabularx}{\textwidth}{l c c c c X}
\toprule
Metric row & Precision & Recall & mAP@0.5 & mAP@0.5:0.95 & Notes/source \\
\midrule
Best observed values & 0.8067 & 0.7401 & 0.7905 & 0.4558 &
Best values reported across the provided YOLO training log summary. \\
Final epoch values & 0.8004 & 0.7298 & 0.7806 & 0.4492 &
Final row from the provided YOLO training log summary. These values describe
object-detection performance, not classification accuracy. \\
\bottomrule
\end{tabularx}
\end{table}


\begin{figure}[htbp]
  \centering
  \includegraphics[width=0.95\textwidth]{yolov8_training_results_curves.png}
  \caption{YOLOv8 training and validation curves for loss and detection metrics.}
  \label{fig:yolo-training-curves}
\end{figure}

Figure~\ref{fig:yolo-training-curves} provides the main training evidence for
the available detector run. The uploaded metric summary reports best observed
precision \(=0.8067\), best recall \(=0.7401\), best mAP@0.5 \(=0.7905\), and
best mAP@0.5:0.95 \(=0.4558\). The final epoch values were precision
\(=0.8004\), recall \(=0.7298\), mAP@0.5 \(=0.7806\), and mAP@0.5:0.95
\(=0.4492\).

\begin{figure}[htbp]
  \centering
  \includegraphics[width=0.80\textwidth]{yolov8_precision_recall_curve.png}
  \caption{Precision--recall curve for fire and smoke detection.}
  \label{fig:yolo-pr-curve}
\end{figure}

\begin{figure}[htbp]
  \centering
  \includegraphics[width=0.80\textwidth]{yolov8_f1_confidence_curve.png}
  \caption{F1-confidence curve used to inspect detection threshold behaviour.}
  \label{fig:yolo-f1-confidence}
\end{figure}

Figures~\ref{fig:yolo-pr-curve} and~\ref{fig:yolo-f1-confidence} help explain
threshold behaviour. Lower thresholds can improve recall but may increase
false positives, while higher thresholds may reduce false alarms but miss
weaker smoke or fire evidence.

\begin{figure}[htbp]
  \centering
  \includegraphics[width=0.78\textwidth]{yolov8_normalized_confusion_matrix.png}
  \caption{Normalized confusion matrix for fire, smoke, and background predictions.}
  \label{fig:yolo-normalized-confusion}
\end{figure}

\begin{figure}[htbp]
  \centering
  \includegraphics[width=0.95\textwidth]{yolov8_labels_vs_predictions_comparison.jpg}
  \caption{Example validation images comparing ground-truth labels with YOLOv8 predictions.}
  \label{fig:yolo-labels-vs-predictions}
\end{figure}

The normalized confusion matrix in Figure~\ref{fig:yolo-normalized-confusion}
and the qualitative examples in Figure~\ref{fig:yolo-labels-vs-predictions}
show how the detector behaves on the available evaluation material. These
figures support the prototype detection workflow, but they do not prove field
performance under all lighting, weather, smoke-density, and camera-motion
conditions. Formal field validation remains future work.

\FloatBarrier

\section{Dashboard and System Validation Results}

The dashboard demonstrates how the FireWatch workflow can be used as an
operator-facing decision-support interface. It brings together risk estimates,
monitoring feeds, detection alerts, system/model status, and map context
instead of leaving each component as a separate tool.

The main validation result from the software layer is integration. Backend
tests verify API behaviour, inference, decision logic, scheduler behaviour,
monitoring endpoints, detection alert persistence, read/unread state, soft
delete behaviour, and separation between detection alerts and prediction
history. A local frontend lint and production build check also confirmed that
the React/TypeScript dashboard compiles successfully.

\begin{table}[htbp]
\centering
\caption{System validation evidence from the repository and local verification.}
\label{tab:system-validation}
\begin{tabularx}{\textwidth}{l X X X}
\toprule
Component & Validation method & Confirmed behaviour & Limitation \\
\midrule
Backend API and models &
Automated backend test suite and smoke check &
Prediction models load, API routes respond, and smoke check confirms key
runtime artefacts. &
Tests are local software checks, not field certification. \\
Scheduler and run persistence &
Backend tests for scheduler configuration and scheduled runs &
Scheduled risk checks use the same persistence path as manual risk checks. &
Long-term unattended operation was not evaluated. \\
SQLite persistence &
Database schema inspection and endpoint tests &
\texttt{run\_history}, \texttt{system\_state}, and
\texttt{detection\_alerts} preserve operational state. &
SQLite is suitable for a local prototype, not multi-user production. \\
Detection alert workflow &
Monitoring endpoint tests &
Alerts support source filtering, latest lookup, read/unread state, and soft
deletion without writing prediction history. &
Visual detection robustness under field lighting remains future work. \\
Dashboard integration &
Frontend lint and production build verification &
The React/TypeScript dashboard builds and displays the integrated monitoring
interface. &
Usability testing with forestry operators was not conducted. \\
\bottomrule
\end{tabularx}
\end{table}


Table~\ref{tab:system-validation} summarizes the system validation evidence.
These checks demonstrate that the academic prototype works as a connected
software system. They do not certify official emergency deployment, field
camera reliability, or fully autonomous drone patrol.

\section{Limitations}

The Stage 1 regression model still has prediction error. Although the
out-of-fold and holdout results show useful performance, unusual weather
patterns and difficult high-FWI days can still produce larger errors. The
predicted FWI value should therefore be interpreted as a decision-support
estimate, not as a replacement for official fire-danger products.

The Stage 2 classifier improves high-risk decision support but does not remove
false positives or false negatives. False positives can increase monitoring
load, while false negatives are more serious because they may reduce attention
during dangerous conditions. This is why FireWatch shows the probability and
decision reason rather than hiding the result behind a single unexplained
label.

The YOLOv8 detection stage is also limited. The uploaded training/evaluation
curves provide useful object-detection evidence, but they are not a substitute
for extended field validation in Karabük forests, under drone motion, varied
lighting, glare, fog, and different smoke densities. Multi-frame confirmation,
glare filtering, bounding-box area filtering, and larger local evaluation
sets remain future work.

The persistence and drone-ready layers are intentionally prototype-oriented.
SQLite is appropriate for a local single-machine demonstration, but a deployed
multi-user system would require a server database, access control, backups,
and operational monitoring. The drone layer supports mock and Tello-oriented
workflows with operator safety gates; it does not claim certified autonomous
drone patrol.

\section{Discussion}

The results answer the research question by showing that FireWatch can combine
daily wildfire risk prediction, visual fire/smoke detection, persistent alert
management, and dashboard-based decision support for the Karabük region. Each
module contributes a different type of evidence. Stage 1 provides continuous
fire-weather risk estimation, Stage 2 supports high-risk classification,
YOLOv8 provides visual evidence, SQLite preserves operational history, and the
dashboard makes the workflow reviewable.

The two-stage prediction design is useful because it separates risk intensity
from high-risk decision support. A continuous FWI estimate helps operators
understand severity, while the Stage 2 classifier focuses on reducing missed
high-risk days near the decision boundary. This is more informative than a
single binary model because it preserves both the estimated risk level and the
operational decision.

The detection model adds a complementary role. Weather-based prediction cannot
confirm that fire or smoke is already visible, and camera detection alone does
not explain whether the wider weather conditions are dangerous. By keeping
prediction and detection separate but visible in the same dashboard, FireWatch
supports both preparedness and visual evidence review.

Overall, FireWatch is strongest as a regional academic decision-support
prototype. It does not claim official emergency integration, field
certification, or production autonomous drone operation. However, the results
demonstrate a complete engineering workflow that connects machine learning,
computer vision, persistence, scheduling, and dashboard monitoring into a
coherent FireWatch prototype.
